\documentclass[12pt]{article}
\usepackage[utf8]{inputenc}
\usepackage[T1]{fontenc}
\usepackage{hyperref}
\title{Ute-Southern Paiute Language Resources for K-12 Curriculum Development}
\author{Compiled by Claude Sonnet 5 (claude-sonnet-5) for Andr\'as Kornai}
\date{}
\begin{document}
\maketitle
\section{Introduction}
Ute-Southern Paiute (ISO 639-3 \texttt{ute}) is a Uto-Aztecan language of the United States. My searches concentrated on Southern Ute and the southern Ute dialects (Southern Ute, Ute Mountain, Uncompahgre). I did not search for Southern Paiute or for Ute Mountain Ute and Northern Ute programmes separately.
\section{Primary Educational Resources}
\begin{itemize}
\item Ute Dictionary, revised edition (2016), by T. Givón with elders Pearl Casías, Mary Inez Cloud and Vida Baker Peabody, published by John Benjamins. It grew out of the Southern Ute Tribe's Ute Language Program, and the Tribe holds the rights to the text. It is reported as one of the main texts used in Ute classes. Preface: \url{https://www.degruyterbrill.com/document/doi/10.1075/clu.15.preface/pdf}
\item Southern Ute Drum report on the new dictionary: \url{https://www.sudrum.com/culture/2016/06/10/new-ute-dictionary-available/}
\item Southern Ute Drum report on the SILDI programme, in which students taught Ute classes and used the dictionary as a reference: \url{https://www.sudrum.com/top-stories/2023/05/19/sildi-program-aims-to-preserve-the-nuu-apaha-pi/}
\end{itemize}
\section{Linguistic and Reference Materials}
\begin{itemize}
\item A Short Dictionary of the Southern Ute Language (J. A. Goss, Southern Ute Tribe, 1961), catalogued in the UBC knowledge base: \url{https://knowledgebase.arts.ubc.ca/?p=3541}
\item UBC wiki page on the Ute Dictionary: \url{https://wiki.ubc.ca/Documentation:RelLex/Ute_Dictionary}
\end{itemize}
\section{Religious and Community Resources}
\begin{itemize}
\item Southern Ute Drum, the tribal newspaper, which covers language and culture: \url{https://www.sudrum.com/?p=11626}
\end{itemize}
\section{Digital and Audio Resources}
I found no audio or video resources in the search results.
\section{Curriculum Development Context}
Ute language classes date back to about 1975, when the Southern Ute Tribe asked Fort Lewis College to host them. The tribe's language programme and the dictionary are the natural base for a curriculum. I did not verify whether current school-level materials exist, and the Southern Ute Tribe's language programme would need to be contacted directly.
\end{document}